\chapter*{General Conclusion}
\addcontentsline{toc}{chapter}{General Conclusion}
\markboth{General Conclusion}{General Conclusion}

This project set out to answer a single practical question: can a legal \ac{AI}
assistant be both deeply grounded in Tunisian law and deployed entirely within a
secure, self-contained local infrastructure? The work carried out over the four
chapters of this report demonstrates that the answer is yes.

Starting from a raw collection of official \ac{JORT} publications, we built a
seven-phase preparation pipeline that extracts, structures, validates, and indexes
Tunisian legal articles into a hybrid vector database. At query time, a compact
Python pipeline runs the user's question through an embedding model, a hybrid
search against the indexed corpus, a cross-encoder reranker, a confidence gate,
and a fine-tuned language model, all without a single external network call. Two
candidate models, Gemma~4~E4B and Qwen3.5~9B, were fine-tuned on a curated
dataset of 1,619 Tunisian legal \ac{QA} examples and served locally in quantized
\ac{GGUF} format via Ollama.

The evaluation confirmed that all three objectives were met. Both models produce
grounded, correctly cited, and disclaimed legal answers in French and Arabic.
The retrieval layer correctly identified the relevant article for every factual
query and declined all out-of-scope questions without hallucinating. Every
component active at inference time runs locally with no external dependencies.
An expert evaluation by a legal professional further confirmed strong performance:
Gemma~4~E4B reached 80\% to 95\% accuracy across direct \ac{QA}, multi-article
synthesis, principle-and-exception, clarification, and in-domain refusal
categories; Qwen3.5~9B held at 80\% on those same categories despite operating
under heavier quantization constraints.

Two limitations remain open for the next development iteration. The indexed corpus
is currently limited to two \ac{JORT} issues, a constraint imposed by the cost of
the GPT-4.1 enrichment step; scaling to a full national legislation index is the
highest-priority next step. Routing is the only category where both models fall
below the 80\% acceptance threshold, Gemma~4~E4B at 73\% and Qwen3.5~9B at 60\%,
which reflects the limited number of routing examples in the training dataset and
points directly to where the next data collection effort should focus.

More broadly, this project demonstrates that sovereign, domain-specific legal
\ac{AI} is technically achievable at the scale of a startup. The combination of
parameter-efficient fine-tuning, hybrid retrieval, and quantized local inference
makes it possible to build a system that is simultaneously specialized, accurate,
and fully self-contained, a meaningful step toward making Tunisian legal knowledge
accessible, reliable, and private.
